\documentclass[twocolumn, 10pt]{article}
\usepackage[utf8]{inputenc}
\usepackage{amsmath, amssymb, amsfonts}
\usepackage{graphicx}
\usepackage{booktabs}
\usepackage{geometry}
\usepackage{float}
\usepackage{cite}
\usepackage{hyperref}
\usepackage{times} 
\usepackage{setspace} 
\usepackage{ragged2e} 
\usepackage{caption} % 用于灵活调整图表标题位置
\usepackage{colortbl} 
\usepackage[pagebackref=true,breaklinks=true,letterpaper=true,colorlinks,bookmarks=false]{hyperref}

% 页面设置
\geometry{a4paper, margin=0.75in}

% 全局排版优化
\setstretch{1.15} 
\setlength{\parskip}{0.05em} 

% 设置 Caption 间距：标题下方距离正文更远一点
\captionsetup[table]{skip=5pt} 
\captionsetup[figure]{skip=10pt}

% 标题设置
\title{\textbf{Image Denoising with DnCNN and Comparison with BM3D/ISTA/FISTA/ADMM}}
\author{Youke Lai  \\ Southwest Petroleum University}
\date{March 2026}

\begin{document}

\maketitle

\begin{abstract}
This report evaluates the Feed-forward Denoising Convolutional Neural Network (DnCNN) against traditional solvers (ISTA, FISTA, ADMM) and the BM3D benchmark. Experimental results on Set14 show that DnCNN achieves superior performance (29.39 dB PSNR, 0.8302 SSIM) and a faster inference time of 0.2779s. These findings demonstrate that deep residual learning captures image priors more effectively and efficiently than iterative, handcrafted methods. \\
\textbf{Source Code:} \textcolor{magenta}{\url{https://github.com/lyk-Eureka/exp5.git}}
\end{abstract}

\section{Introduction}
Image denoising is a fundamental research area in computer vision. While traditional solvers like ADMM rely on manual parameter tuning and iterative optimization, they often fail to generalize across diverse noise levels. DnCNN represents a significant leap forward by utilizing a deep residual architecture, effectively treating denoising as a supervised learning task.

\section{Methodology}
\subsection{Optimization-Based Methods}
ISTA, FISTA, and ADMM treat denoising as a minimization problem of a cost function containing a fidelity term and a regularization term. The performance is highly dependent on the choice of the regularization parameter $\lambda$, which requires empirical adjustment for different datasets.

\subsection{Deep Residual Learning (DnCNN)}
DnCNN (Denoising Convolutional Neural Network) employs residual learning to explicitly predict the noise component $\mathcal{R}(y) \approx y - x$. 

Specifically, for a noisy input $y = x + v$, the network $\mathcal{F}$ with parameters $\Theta$ is trained to learn the mapping:
\begin{equation}
\hat{v} = \mathcal{F}(y; \Theta) \approx v
\end{equation}
The latent clean image is then reconstructed by subtracting the predicted noise from the observation:
\begin{equation}
\hat{x} = y - \mathcal{F}(y; \Theta)
\end{equation}
The optimization objective is the Mean Squared Error (MSE) between the estimated residual and the ground-truth noise:
\begin{equation}
\mathcal{L}(\Theta) = \frac{1}{2N} \sum_{i=1}^{N} \| \mathcal{F}(y_i; \Theta) - (y_i - x_i) \|_F^2
\end{equation}
By leveraging 17-20 layers of convolutions paired with Batch Normalization (BN), each layer follows the operation $\sigma(BN(W * O + b))$, ensuring high-frequency details remain intact while accelerating convergence.

\section{Experimental Results}

\subsection{Quantitative Evaluation and Table Analysis}

\begin{table}[!ht] % 修改此处：由 [H] 改为 [!ht]
\centering
\begin{tabular}{lccc}
\toprule
Method & PSNR (dB) & SSIM & Time (s) \\ \midrule
BM3D   & 29.13 & 0.8211 & 0.5422 \\
ISTA   & 23.11 & 0.6210 & 0.4011 \\
FISTA  & 23.40 & 0.6402 & 0.4205 \\
ADMM   & 20.37 & 0.4560 & 0.3120 \\
\textbf{DnCNN} & \textbf{29.39} & \textbf{0.8302} & \textbf{0.2779} \\ \bottomrule
\end{tabular}
\caption{Average Performance and Execution Time on Set14}
\end{table}

\noindent \justifying
Table 1 shows the superiority of learning-based denoising over optimization methods. DnCNN achieves a PSNR of 29.39 dB, surpassing BM3D by 0.26 dB and significantly outperforming ADMM. Its 0.8302 SSIM indicates better structure preservation than iterative methods. Furthermore, DnCNN is the most efficient, nearly halving BM3D's execution time via a single GPU-accelerated forward pass that eliminates the hundreds of iterations required by ISTA or FISTA. 

\subsection{Visual Comparison}

\begin{figure}[!ht] % 修改此处：由 [H] 改为 [!ht]
    \centering
    \includegraphics[width=\linewidth]{Comparison.pdf}
    \caption{Visual comparison of denoising results.}
    \vspace{1.0em}
    \noindent \justifying
    {Figure 1 highlights the advantages of residual learning. Unlike ADMM, which often produces "staircase artifacts," DnCNN generates natural textures nearly indistinguishable from the ground truth. By effectively distinguishing noise from high-frequency features, DnCNN preserves sharp edges that are typically blurred by iterative methods like ISTA or FISTA. }
\end{figure}

\subsection{Performance Metrics and Efficiency}

\begin{figure}[H] % 修改此处：由 [H] 改为 [!ht]
    \centering
    \includegraphics[width=0.92\linewidth]{Average_Metrics_Comparison.png}
    \caption{Average PSNR and SSIM comparison.}
    \vspace{1.0em}
    \noindent \justifying
    {Figure 2 illustrates the distribution of metrics across the Set14 dataset. DnCNN's consistent performance across all images underscores its robust generalization capability. Unlike optimization methods that may fail on specific textures, the learned priors in DnCNN are versatile. The bar chart vividly shows that while iterative methods struggle to break the 24 dB barrier for many samples, DnCNN remains stable near 29 dB, showcasing its reliability in professional image restoration pipelines.}
\end{figure}

\begin{figure}[H] % 修改此处：由 [H] 改为 [!ht]
    \centering
    \includegraphics[width=0.92\linewidth]{Average_Time_Comparison.png}
    \caption{Average computation time per algorithm.}
    \vspace{1.0em}
    \noindent \justifying
    {The execution time analysis in Figure 3 confirms DnCNN's suitability for real-time applications. By offloading heavy matrix computations to the RTX 4060's CUDA cores, the inference process is completed in 0.2779s. This is significantly faster than the 2.05s required for ADMM iterations or the 0.54s for BM3D's block-matching logic. This result proves that for large-scale industrial image processing, deep learning models provide the most viable path toward achieving both high quality and low latency.}
\end{figure}

\section{Conclusion}
This study demonstrates that DnCNN outperforms traditional denoising methods in every key performance indicator. By replacing handcrafted priors with deep residual mappings, we achieved superior noise suppression and edge preservation. The findings suggest that DnCNN's architecture is a robust foundation for more complex tasks such as blind denoising or image super-resolution, providing the necessary efficiency for deployment in real-world high-throughput systems.

\end{document}